\subsection{The two marginal notes on 57:310: mathematics and reading}

The multiplication margin is diplomatically transcribed as
``La multiplication par un multiple de $p$ est [surjective?]
[two illegible words] [ouverte?] (?) [comp. irr\'eductibles?]''.
The original image confirms ``Non, exemple'' and the branch diagram,
but does not justify removing the marked uncertainties. The mathematical
universal-openness question is nevertheless answered negatively by
Theorem E4.H, including for a smooth projective family. Its status is
therefore separate from the unresolved transcription of the qualifier.

The bracketed finite-type criterion has the alternatives ``au dessus''
and ``en dehors'' of an open set, with an uncertain dense/constructible
qualification. In the proper-flat-finitely-presented,
cohomologically-flat-in-degree-zero Picard setting over a noetherian
base, finite presentation of $\operatorname{Pic}^{\tau}$ already proves
that every such base change, open restriction, and closed complement has
finitely many irreducible components. Thus either indicated
finite-component condition holds in this Picard setting. This does not
prove a proposed general criterion for arbitrary group schemes.

\paragraph{An explicit warning about the abstract generalization.}
Over $S=\mathbf A^1_{\mathbf Q}$, enumerate the distinct points
$s_i=(t-i)$, $i\geq1$. Glue a copy $S_0$ of $S$ and copies $S_i$,
$i\geq1$, by identifying $S_i\setminus\{s_i\}$ with the corresponding
open of $S_0$. Call the resulting scheme $G$. Its fibre over $s_i$ has
two points, and every other fibre has one. It is an \'etale commutative
group scheme with the two-point fibre given the group law
$\mathbf Z/2$: on two different charts $S_i,S_j$, their sum is the
section that uses the extra origin at $s_i$ and $s_j$, obtained by gluing
on $S\setminus\{s_j\}$ and $S\setminus\{s_i\}$. On the same chart
the sum is the zero section. These formulas agree on chart intersections;
the group laws hold fibrewise and hence hold as morphisms into the
\'etale space.

The scheme $G$ is irreducible, since all its charts have the same dense
generic point, but is not quasi-compact: any finite collection of these
charts misses infinitely many extra origins. For every closed irreducible
$S'\subset S$, its restriction over a nonempty open of $S'$ has finitely
many irreducible components. Consequently this generic-open
component-count condition does not characterize finite type for arbitrary
locally finite-type group schemes. This is deliberately not presented as
a Picard realization. It only rules out silently replacing the uncertain
margin by that unrestricted abstract assertion.
